\chapimg{fig/op12.png}
\IfFileExists{extra/lead12.tex}{\chapterlead{\subhd{After this chapter you can} design a drip system from the emitter outward: relate discharge to pressure, compute lateral head loss with Hazen--Williams and the Christiansen factor, select the lateral diameter by emission uniformity, compute the set time, and size the manifold and filter.}
}{}
\chapter{Drip Irrigation Design}\label{chap:12}
\markboth{Drip Irrigation Design}{Drip Irrigation Design}
\IfFileExists{extra/bg12.tex}{\section{Introduction}
Drip (trickle) irrigation delivers water at a few litres per hour to the root zone through emitters on small tubes. The design follows the water from the emitter backwards: the emitter sets the discharge, the lateral the pressure variation, the submain and main the pressure at the head, and the filter and pump the energy. A design is good when the emission uniformity is above about 90 per cent.

\section{Types of drip systems}
Drip is of three kinds: \emph{point-source} emitters (drippers on the line or in-line), \emph{line-source} (dripperline with continuous emitters) and \emph{subsurface} (SDI, buried). Emitters may be \emph{pressure compensating} (the discharge is constant over a range of pressures, $x\approx0$) or \emph{non-compensating} ($x$ from 0.5 to 1).

\section{Emitter hydraulics}
The discharge of an emitter is $q=K_dH^{x}$. A turbulent-flow emitter has $x=0.5$: with $q_0=1.6$ L\,h$^{-1}$ at $H_0=10$ m, a loss of 1.5 m of head gives $q=1.48$ L\,h$^{-1}$, a fall of 7.8 per cent.

\section{Lateral design}
The Hazen--Williams head loss in a pipe is $h_f=10.67\,L\,Q^{1.852}/(C^{1.852}D^{4.871})$ with $Q$ in m$^3$\,s$^{-1}$ and $D$ in m. For a lateral with outlets the loss is multiplied by the Christiansen factor $F=\frac{1}{m+1}+\frac{1}{2N}+\frac{\sqrt{m-1}}{6N^2}$, which for $m=1.852$ and $N=200$ outlets is 0.353. The lateral of 100 m with 200 emitters carries 0.089 L\,s$^{-1}$.

\section{Emission uniformity}
The \kw{emission uniformity} is $EU=100(1-1.27\,CV/\sqrt{n_e})\,q_{min}/q_{avg}$. It combines the manufacturing variation $CV$ (here 0.03) and the pressure variation along the lateral. An EU above 90 per cent is excellent, 80--90 acceptable and below 70 unacceptable.

}{}
\IfFileExists{extra/deep12.tex}{\section{Drip irrigation: what the layer does}
The drip layer designs the system from the emitter to the root zone. From the soil of the field it takes the water that must be applied (the net depth $Z_n$ to the trigger), the water-holding properties ($\theta_{33}$, $\theta_{1500}$, $K_s$) and the root depth; from the user it takes the emitter, the spacing of emitters and laterals, the length and diameter of the lateral, the inlet pressure and the depth of placement. It returns the emission uniformity, the head lost along the lateral, the wetted bulb under the average emitter, the run time, and the fate of every litre: stored in the root zone, deep percolation, or evaporation from the wet surface.

\section{Types of emitters and their parameters}
The application offers six emitters. They differ in the variation of the discharge from one emitter to another (the manufacturing coefficient of variation $c_v$), in the sensitivity of the discharge to pressure (the exponent $x$), in the number of emitters that a plant is watered by ($n_e$) and, for the pressure-compensating drippers, in the minimum head $H_{\min}$ below which compensation fails.

\begin{table}[!htbp]\centering\footnotesize\setlength{\tabcolsep}{3pt}\caption{Emitters of the drip layer.}
\begin{tabular}{lrrrrrrr}\toprule
\hd{Emitter} & \hd{$c_v$} & \hd{$x$} & \hd{$n_e$} & \hd{$H_{\min}$ (m)} & \hd{default $q$ (L/h)} & \hd{default $S_e$ (m)} & \hd{inside $D$ (mm)}\\\midrule
Online emitter & 0.07 & 0.50 & 1 & 0.0 & 4.0 & 0.75 & 16.0\\
Inline emitter & 0.04 & 0.50 & 1 & 0.0 & 2.0 & 0.50 & 16.0\\
Drip tape & 0.08 & 0.60 & 2 & 0.0 & 1.0 & 0.20 & 16.0\\
PC AS dripper & 0.03 & 0.00 & 1 & 5.1 & 2.0 & 0.50 & 13.7\\
PC LCNL dripper & 0.03 & 0.00 & 1 & 10.2 & 2.0 & 0.50 & 13.7\\
PC HCNL dripper & 0.03 & 0.00 & 1 & 14.3 & 3.0 & 0.50 & 13.7\\\bottomrule\end{tabular}\end{table}

\textbf{Online emitters} are punched into a polyethylene tube; they are cheap and can be placed where the plants are, but they vary most ($c_v=0.07$). \textbf{Inline emitters} are welded in the tube at a fixed spacing and vary less ($c_v=0.04$). \textbf{Drip tape} is a thin-walled tube with a continuous labyrinth; it is used for one season on vegetables at close spacing (0.2 m), and because the wall is thin the discharge is sensitive to pressure ($x=0.6$) and each plant is watered by two emitters. The three \textbf{pressure-compensating} drippers (PC AS, anti-siphon; PC LCNL, low-pressure non-leakage; PC HCNL, high-pressure non-leakage) keep the discharge constant ($x=0$) above a minimum head, which is why they are used on slopes and in long laterals. The non-leakage types shut off below 0.15 or 0.3 m of head, which prevents draining of the lateral at the end of the run. The lateral pipes of the library are 16 mm (inside 13.7 mm, wall 0.95 mm, working pressure 4.5 bar) and 20 mm (inside 17.5 mm).

\section{Where to place the emitter}
The emitter can lie on the surface or be buried at 15, 25, 45 or 55 cm. Burying reduces the evaporation from the wet surface and keeps the water in the root zone; but if the emitter is deeper than the roots, or the soil is coarse, the water leaves the root zone before the roots can take it. The application computes, for every placement, the water that passes below the maximum rooting depth ($R_{max}$) as deep percolation, and the water that reaches the surface as evaporation.

\section{Governing equations}
\textbf{The emitter.} The discharge depends on the head $H$ at the emitter as a power law,

\begin{equation}\boxed{q=q_0\left(\frac{H}{H_0}\right)^{x},}\label{eq:d-em}\end{equation}
\par\noindent where $q$ is the discharge of the emitter (L\,h$^{-1}$), $q_0$ the nominal discharge at the nominal head $H_0$ (m), $H$ the head at the emitter (m) and $x$ the emitter exponent ($x=0.5$ for a fully turbulent orifice, 1 for a laminar one, 0 for a perfect compensator). A pressure-compensating emitter keeps $q=q_0[1+0.04(H-H_0)/H_0]$ for $H\ge H_{\min}$ and falls as $q_0(H/H_{\min})^{0.5}$ below it.



\textbf{The lateral.} The head falls along the lateral because of friction. For the segment between emitters $j$ and $j+1$, of length $S_e$ and carrying the flow still to be delivered downstream, $Q_j=\sum_{i>j}q_i$, the Hazen--Williams loss is

\begin{equation}\boxed{\Delta H_j=\frac{10.67\,S_e\,Q_j^{1.852}}{C^{1.852}D^{4.871}},\qquad H_{j+1}=H_j-\Delta H_j,}\label{eq:d-hw}\end{equation}
\par\noindent where $\Delta H_j$ is the head lost in the segment (m), $S_e$ the emitter spacing (m), $Q_j$ the flow in the segment (m$^3$\,s$^{-1}$; the sum of the discharges, in L\,h$^{-1}$, of the emitters beyond $j$ divided by $3.6\times10^{6}$), $C$ the Hazen--Williams coefficient (140 in the application) and $D$ the inside diameter (m).



The discharges depend on the heads and the heads on the discharges, so the application iterates (four undamped passes, then under-relaxed with a factor 0.5, until the discharges change by less than $10^{-7}q$, at most 500 passes). The result is the head and the discharge at every emitter. A lateral is outside the usual design limit when the head lost exceeds 20 per cent of the inlet head: $h_f/H_0>0.2$.

\textbf{Uniformity.} The coefficient of variation of the discharge combines the manufacturing variation and the variation due to pressure,

\begin{equation}\sigma_q=\sqrt{c_v^2+\sigma_P^2},\quad\sigma_P^2=\frac{1}{N+1}\sum_j\left(\frac{q_j}{\bar q}-1\right)^2,\quad E_u=\min\left\{1,\ \left(1-\frac{1.27\,c_v}{\sqrt{n_e}}\right)\frac{q_{\min}}{\bar q}\right\},\label{eq:d-eu}\end{equation}
\par\noindent where $\sigma_q$ is the total coefficient of variation of the discharge, $c_v$ the manufacturing coefficient, $\sigma_P$ that due to the pressure, $q_j$ the discharge of emitter $j$, $\bar q$ the mean discharge (L\,h$^{-1}$), $N+1$ the number of emitters on the lateral, $E_u$ the emission uniformity (limited to the range 0.3 to 1), $n_e$ the number of emitters per plant and $q_{\min}$ the smallest discharge on the lateral.



The lateral carries $Q=N\bar q/3600$ L\,s$^{-1}$.

\begin{table}[!htbp]\centering\footnotesize\setlength{\tabcolsep}{3pt}\caption{Laterals with emitters of 2 L/h at 0.5 m spacing and 10 m inlet head: head loss and uniformity from the iteration.}
\begin{tabular}{lrrrrrrr}\toprule
\hd{Emitter} & \hd{$L$ (m)} & \hd{$D$ (mm)} & \hd{$h_f$ (m)} & \hd{$h_f/H_0$ (\%)} & \hd{$q_{\min}$ (L/h)} & \hd{$\bar q$ (L/h)} & \hd{$E_u$ (\%)}\\\midrule
Inline emitter & 100 & 13.7 & 1.91 & 19.1 & 1.799 & 1.852 & 92.2\\
Inline emitter & 200 & 13.7 & 6.82 & 68.2 & 1.128 & 1.366 & 78.4\\
PC AS dripper & 200 & 13.7 & 8.26 & 82.6 & 1.168 & 1.531 & 73.4\\
Inline emitter & 100 & 16 & 0.98 & 9.8 & 1.900 & 1.926 & 93.6\\
Inline emitter & 200 & 16 & 4.67 & 46.7 & 1.460 & 1.605 & 86.3\\
PC AS dripper & 200 & 16 & 6.32 & 63.2 & 1.698 & 1.856 & 88.0\\
Inline emitter & 100 & 17.5 & 0.65 & 6.5 & 1.934 & 1.951 & 94.1\\
Inline emitter & 200 & 17.5 & 3.51 & 35.1 & 1.611 & 1.715 & 89.2\\
PC AS dripper & 200 & 17.5 & 4.78 & 47.8 & 1.962 & 1.972 & 95.7\\
Inline emitter & 100 & 20 & 0.35 & 3.5 & 1.965 & 1.974 & 94.5\\
Inline emitter & 200 & 20 & 2.13 & 21.3 & 1.774 & 1.834 & 91.8\\
PC AS dripper & 200 & 20 & 2.53 & 25.3 & 1.980 & 1.985 & 95.9\\\bottomrule\end{tabular}\end{table}

\exbox{Head loss and uniformity of a lateral}{An inline-emitter lateral ($q_0=2$ L\,h$^{-1}$, $S_e=0.5$ m, $L=100$ m, $D=16$ mm, $H_0=10$ m, $x=0.5$, $c_v=0.04$) is run in the application. Report the number of emitters, the head lost, the extreme and mean discharges, $E_u$ and the flow.}\sol \givenline{$q_0=2$ L\,h$^{-1}$ ; $S_e=0.5$ m ; $L=100$ m ; $D=16$ mm ; $H_0=10$ m ; $x=0.5$ ; $c_v=0.04$}\ans $N=L/S_e=200$ intervals, i.e.\ 201 emitters \ans At the inlet the segment carries the discharge of the 200 emitters beyond: $Q\approx200\times2/3.6\times10^6=1.11\times10^{-4}$ m$^3$\,s$^{-1}$, $\Delta H=10.67(0.5)(1.11\times10^{-4})^{1.852}/(140^{1.852}\times0.016^{4.871})=0.0150$ m \ans Adding the losses of all segments with the discharges updated by Eq.~(\ref{eq:d-em}) gives, after convergence, $h_f=0.98$ m, so the last emitter has $H=9.02$ m, i.e.\ $h_f/H_0=9.8$ per cent (within the 20 per cent limit) \ans $q_{\min}=1.900$ L\,h$^{-1}$ (last emitter), $\bar q=1.926$ L\,h$^{-1}$ \ans $\sigma_P$ and the manufacturing variation give $\sigma_q=0.0427$ \ans $E_u=(1-1.27\times0.04/\sqrt1)(1.900/1.926)=0.936$, that is 93.6 per cent \ans Flow $Q=201\times1.926/3600=0.108$ L\,s$^{-1}$ \finans

\section{The wetted bulb}
Around a buried or surface emitter the water spreads as a body of revolution. The application describes it as a half-ellipsoid below the emitter (radius $r$ at the emitter depth, vertical semi-axis $a\,r$) and a short one above it (vertical semi-axis $u\,r$, $u=0.6$). The width of the body at depth $z$ is

\begin{equation}\boxed{\rho(z)=r\sqrt{1-\left(\frac{z-z_d}{a\,r}\right)^2}\ \ (z\ge z_d),\qquad \rho(z)=r\sqrt{1-\left(\frac{z_d-z}{u\,r}\right)^2}\ \ (z<z_d),}\label{eq:d-bulb}\end{equation}
\par\noindent where $\rho(z)$ is the radius of the wet body at depth $z$ (mm), $z_d$ the depth of the emitter (mm), $r$ the radius at emitter depth (mm), $a$ the aspect ratio (depth below the emitter over radius) and $u=0.6$ the same ratio above.



The volume of soil that holds the water of one emitter is $V_w=V/\Delta\theta$, where $V$ is the volume applied and $\Delta\theta=\max\{0.03,\ f_T(\theta_{33}-\theta_{1500})\}$ is the water that the soil takes up between the trigger and field capacity; the radius $r$ is the one at which the volume of the body, $\int\pi\rho^2dz$ (or $\int2\rho\,dz$ for a line source), equals $V_w$. The shape depends on the soil through the aspect ratio, which expresses the competition of gravity and capillarity:

\begin{equation}\boxed{a=\mathrm{clip}\left\{\left[0.7+1.6\left(\frac{K_s}{K_s+15}\right)^{1.1}\right]\sqrt{\mathrm{clip}\!\left(\frac{2}{q};\,0.7,\,1.4\right)};\ 0.75,\,2.4\right\},}\label{eq:d-asp}\end{equation}
\par\noindent where $a$ is the aspect ratio of the bulb, $K_s$ the saturated conductivity in cm\,d$^{-1}$ ($=2.4\times K_s$ in mm\,h$^{-1}$) and $q$ the discharge of the emitter (L\,h$^{-1}$); coarse soils give a narrow, deep bulb ($a$ large), fine soils a wide, shallow one, and a larger discharge spreads the water more.



When the emitter spacing is less than 0.34 m, or the bulbs of neighbouring emitters touch ($2r>S_e$ with $S_e<1.2$ m), the emitters form a continuous wet strip and the body is treated as a line source with the volume per metre $V/S_e$.

\textbf{Fate of the water.} Of the volume $V_i$ applied by one emitter, the part of the body below the maximum rooting depth $R_m$ is deep percolation, a part $e=\min\{0.5V_i,\ A_{s}E_0\cdot2\}$ evaporates from the wet surface of area $A_s$ ($E_0$ is the surface evaporation rate in mm\,d$^{-1}$), and the rest, $V_i-V_{below}-e$, is the water in the root zone. The part of that held within the requirement $V_{req}=Z_nS_eS_l$ is \emph{stored}; the excess is deep percolation too. The emitters are not equal, so five emitters at the quantiles $z=-1.28,-0.52,0,0.52,1.28$ of the normal distribution, with discharge $V(1+\sigma_qz)$, are computed and averaged:

\begin{equation}E_a=\frac{\overline{V_{st}}}{\overline V},\quad E_s=\frac{\overline{V_{st}}}{V_{req}},\quad DP=\frac{\overline{V_{below}+V_{excess}}}{\overline V},\quad t=\frac{V_{req}}{\bar q\,E_u},\label{eq:d-fate}\end{equation}
\par\noindent where $E_a$ is the application efficiency, $E_s$ the storage efficiency, $DP$ the deep-percolation share, $V_{st}$ the volume stored in the root zone of one quantile emitter (L), $V$ the volume it delivers (L), $V_{req}=Z_nS_eS_l$ the requirement per emitter (L, with $Z_n$ in mm and spacings in m), the overbar the average over the five quantiles and $t$ the default run time (h); the factor $E_u$ lengthens the run so that the poorest emitter still refills the root zone.



\begin{table}[!htbp]\centering\footnotesize\setlength{\tabcolsep}{3pt}\caption{Drip design on the ten soils: inline emitter 2 L/h at 0.5 m, laterals every 1.5 m, buried 25 cm, net depth 9 mm.}
\begin{tabular}{llrrrrrrrr}\toprule
\hd{ID} & \hd{Soil} & \hd{$K_s$ (cm/d)} & \hd{$a$} & \hd{run $t$ (h)} & \hd{width (cm)} & \hd{depth (cm)} & \hd{$E_a$ (\%)} & \hd{$DP$ (\%)} & \hd{Evap. (\%)}\\\midrule
S01 & Coarse loam & 31.6 & 1.74 & 3.7 & 50 & 68 & 92 & 8 & 0\\
S02 & Coarse sand & 275.4 & 2.21 & 3.7 & 66 & 98 & 68 & 32 & 0\\
S03 & Sand & 323.3 & 2.22 & 3.7 & 66 & 98 & 67 & 33 & 0\\
S04 & Loamy sand & 207.3 & 2.18 & 3.7 & 66 & 97 & 68 & 32 & 0\\
S05 & Sandy loam & 107.3 & 2.09 & 3.7 & 54 & 81 & 79 & 21 & 0\\
S06 & Loam & 29.0 & 1.71 & 3.7 & 48 & 66 & 93 & 7 & 0\\
S07 & Silt loam & 26.2 & 1.67 & 3.7 & 41 & 60 & 93 & 7 & 0\\
S08 & Sandy clay loam & 12.7 & 1.38 & 3.7 & 59 & 66 & 93 & 7 & 0\\
S09 & Clay loam & 7.7 & 1.19 & 3.7 & 54 & 57 & 94 & 6 & 0\\
S10 & Silty clay loam & 9.0 & 1.24 & 3.7 & 49 & 56 & 94 & 6 & 0\\\bottomrule\end{tabular}\end{table}

The coarse soils give the narrowest and deepest bulbs, and the largest losses to deep percolation; the fine soils spread sideways and hold the water in the root zone. This is why the placement and the emitter spacing must be chosen soil by soil.

\begin{table}[!htbp]\centering\footnotesize\setlength{\tabcolsep}{3pt}\caption{Effect of the depth of placement on the loam S06.}
\begin{tabular}{rrrrrrr}\toprule
\hd{Placement (cm)} & \hd{run $t$ (h)} & \hd{bulb depth (cm)} & \hd{$E_a$ (\%)} & \hd{$E_s$ (\%)} & \hd{$DP$ (\%)} & \hd{Evaporation (\%)}\\\midrule
surface & 3.7 & 48 & 81 & 86 & 0 & 19\\
15 & 3.7 & 56 & 94 & 100 & 6 & 0\\
25 & 3.7 & 66 & 93 & 99 & 7 & 0\\
45 & 3.7 & 86 & 60 & 64 & 40 & 0\\
55 & 3.7 & 96 & 38 & 40 & 62 & 0\\\bottomrule\end{tabular}\end{table}

\begin{table}[!htbp]\centering\footnotesize\setlength{\tabcolsep}{3pt}\caption{Effect of the emitter spacing on the loam S06 ($S_l=1.5$ m).}
\begin{tabular}{rrrrrr}\toprule
\hd{Emitter spacing (m)} & \hd{run $t$ (h)} & \hd{$E_a$ (\%)} & \hd{$E_s$ (\%)} & \hd{$DP$ (\%)} & \hd{emitters per ha}\\\midrule
0.20 & 1.8 & 88 & 99 & 12 & 33333\\
0.30 & 2.4 & 91 & 100 & 9 & 22222\\
0.50 & 3.7 & 93 & 99 & 7 & 13333\\
0.75 & 5.5 & 88 & 93 & 12 & 8889\\
1.00 & 7.2 & 84 & 88 & 16 & 6667\\
1.50 & 10.7 & 78 & 82 & 22 & 4444\\\bottomrule\end{tabular}\end{table}

\begin{table}[!htbp]\centering\footnotesize\setlength{\tabcolsep}{3pt}\caption{Effect of the emitter discharge on the loam S06.}
\begin{tabular}{rrrrrr}\toprule
\hd{$q$ (L/h)} & \hd{$a$} & \hd{run $t$ (h)} & \hd{width (cm)} & \hd{$E_a$ (\%)} & \hd{$DP$ (\%)}\\\midrule
1.0 & 2.02 & 7.2 & 45 & 90 & 10\\
1.6 & 1.91 & 4.6 & 46 & 91 & 9\\
2.0 & 1.71 & 3.7 & 48 & 93 & 7\\
3.0 & 1.43 & 2.6 & 52 & 92 & 8\\
4.0 & 1.43 & 2.1 & 52 & 90 & 10\\
8.0 & 1.43 & 1.5 & 55 & 80 & 20\\\bottomrule\end{tabular}\end{table}

\exbox{A drip design from soil to run time}{Design a drip system for the loam S06 for a crop that needs $Z_n=9$ mm per irrigation: inline emitters 2 L\,h$^{-1}$ at 0.5 m, laterals 1.5 m apart, 100 m long, 16 mm, 10 m head, buried at 25 cm, root zone 60 cm. Report the run time, the bulb and the fate of the water.}\sol \givenline{$Z_n=9$ mm per irrigation}\ans Requirement per emitter $V_{req}=9\times0.5\times1.5=6.75$ L \ans $E_u=0.936$, $\bar q=1.926$ L\,h$^{-1}$, so $t=V_{req}/(\bar qE_u)=6.75/(1.926\times0.936)=3.74$ h \ans Volume per emitter $V=\bar qt=7.21$ L, i.e.\ a gross depth of $9.6$ mm \ans Aspect ratio of the loam ($K_s=29.0$ cm\,d$^{-1}$): $a=1.71$ \ans The average emitter wets a width of $48$ cm to a depth of $66$ cm \ans $E_a=93.0$, $E_s=99.3$, $DP=7.0$, evaporation $0.0$ per cent \ans Per hectare: $13333$ emitters, $6667$ m of lateral, system flow $7.13$ L\,s$^{-1}$\,ha$^{-1}$ \finans

}{}
\IfFileExists{extra/deep212.tex}{\section{System design: zones, sets and pumping}
The lateral and the emitter fix the discharge per unit area; the farm supply fixes how much of the farm can be watered at once. The application divides the field into zones that are irrigated in turn, and computes the length of the round. Let $A_f$ be the area to irrigate (ha), $Q_s$ the supply (L\,s$^{-1}$), $q_{ha}$ the flow of the design per hectare (L\,s$^{-1}$\,ha$^{-1}$), $t$ the run time of a zone (h) and $T_i$ the irrigation interval (d). Then

\begin{equation}q_{ha}=\frac{10^{4}}{S_eS_l}\,\frac{\bar q}{3600},\qquad N_z=\left\lceil\frac{A_fq_{ha}}{Q_s}\right\rceil,\qquad t_{round}=N_z\,t,\qquad \text{feasible if }t_{round}\le 20\,T_i,\label{eq:d-zone}\end{equation}
\par\noindent where $q_{ha}$ is the flow per hectare (L\,s$^{-1}$\,ha$^{-1}$), $S_e$ and $S_l$ the emitter and lateral spacings (m), $\bar q$ the mean emitter discharge (L\,h$^{-1}$), $A_f$ the area (ha), $Q_s$ the farm supply (L\,s$^{-1}$), $N_z$ the number of zones (sets), $t$ the run time of one zone (h), $t_{round}$ the duration of a round (h), $T_i$ the interval (d) and 20 h the working day of the system.



The seasonal volume follows from the gross depth: $\sum d_g=n_{irr}\,d_n/E_a$ per season, in mm, or $10\sum d_g$ m$^3$ per hectare. The pump power is $P=\rho gQH/\eta$ with $H$ the sum of the lift, the head at the inlet of the lateral, the friction in the submain and mainline, and the filter loss.

\begin{table}[!htbp]\centering\footnotesize\setlength{\tabcolsep}{3pt}\caption{Zones and round for the inline-emitter design of Chapter~13 ($q_{ha}=7.15$ L\,s$^{-1}$\,ha$^{-1}$, run time 3.74 h).}
\begin{tabular}{rrrrrrl}\toprule
\hd{$A_f$ (ha)} & \hd{$Q_s$ (L/s)} & \hd{$q_{ha}$ (L/s/ha)} & \hd{zones} & \hd{round (h)} & \hd{round (d at 20 h)} & \hd{fits $T_i=2$ d}\\\midrule
1 & 10 & 7.15 & 1 & 3.7 & 0.2 & yes\\
1 & 30 & 7.15 & 1 & 3.7 & 0.2 & yes\\
2 & 10 & 7.15 & 2 & 7.5 & 0.4 & yes\\
2 & 30 & 7.15 & 1 & 3.7 & 0.2 & yes\\
5 & 10 & 7.15 & 4 & 15.0 & 0.7 & yes\\
5 & 30 & 7.15 & 2 & 7.5 & 0.4 & yes\\
10 & 10 & 7.15 & 8 & 29.9 & 1.5 & yes\\
10 & 30 & 7.15 & 3 & 11.2 & 0.6 & yes\\
20 & 10 & 7.15 & 15 & 56.1 & 2.8 & no\\
20 & 30 & 7.15 & 5 & 18.7 & 0.9 & yes\\
40 & 10 & 7.15 & 29 & 108.5 & 5.4 & no\\
40 & 30 & 7.15 & 10 & 37.4 & 1.9 & yes\\\bottomrule\end{tabular}\end{table}

\exbox{Zones and pumping}{A farm of 10 ha is to be watered by the drip design of the loam (emitters 0.5 m, laterals 1.5 m, $\bar q=1.93$ L\,h$^{-1}$, run time 3.74 h) from a supply of 30 L\,s$^{-1}$, every 2 days. Find the flow per hectare, the number of zones, the length of the round and the pump power for a total head of 45 m at 70 per cent efficiency.}\sol \givenline{$\bar q=1.93$ L\,h$^{-1}$}\ans \[q_{ha}=10^4/(0.5\times1.5)\times1.93/3600=7.15\] L\,s$^{-1}$\,ha$^{-1}$ ($13333$ emitters per ha) \ans Whole farm: $10\times7.15=71.5$ L\,s$^{-1}$, which is below the supply of 30, so one zone could be run if the pipes allowed: $N_z=\lceil71.5/30\rceil=3$ \ans Round $3\times3.74=11.2$ h, well within $20\times2=40$ h \ans Better to split in zones of equal size to keep pipe sizes small: with 4 zones each flows $17.9$ L\,s$^{-1}$ \ans Pump power \[P=\rho gQH/\eta=1000\times9.81\times0.030\times45/0.70=18.9\] kW \finans

}{}
\IfFileExists{extra/deep312.tex}{\section{Salinity under drip irrigation}
Drip irrigation keeps the soil wet at the emitter and drives the salt to the edge of the wetted bulb, where it accumulates between irrigations and where the roots do not go. At the emitter the soil water has nearly the salinity of the irrigation water, $EC_w$; at the edge of the bulb it may be several times larger. The salt balance of the bulb, in the steady state, is $EC_{dw}=EC_w/LF$ for the drainage water of salinity $EC_{dw}$ at a leaching fraction $LF$, and the average salinity of the root zone is intermediate. The yield response of a crop to salinity is the piecewise-linear law of Maas and Hoffman,

\begin{equation}\boxed{\frac{Y}{Y_m}=\max\left\{0,\ 1-\frac{b}{100}\,(EC_e-a)\right\}\quad(EC_e>a),}\label{eq:d-mh}\end{equation}
\par\noindent where $Y/Y_m$ is the relative yield, $EC_e$ the mean salinity of the saturation extract of the root zone (dS\,m$^{-1}$), $a$ the threshold salinity below which there is no loss (dS\,m$^{-1}$) and $b$ the slope, the percentage loss per dS\,m$^{-1}$ above the threshold.



With the mean root-zone salinity estimated as $EC_e\approx\tfrac32EC_w$ at a leaching fraction of 0.15 to 0.2 (the usual rule of thumb for a leaching requirement met), the table gives the relative yield of crops for three irrigation waters.

\begin{table}[!htbp]\centering\footnotesize\setlength{\tabcolsep}{3pt}\caption{Relative yield from Eq.~\ref{eq:d-mh} with $EC_e=1.5\,EC_w$ (typical published values of $a$ and $b$).}
\begin{tabular}{lrrrrr}\toprule
\hd{Crop} & \hd{threshold $a$} & \hd{slope $b$} & \hd{$Y/Y_m$ (\%) at $EC_w=0.5$} & \hd{1.5} & \hd{3.0}\\\midrule
Maize & 1.7 & 12 & 100 & 93 & 66\\
Wheat & 6.0 & 7.1 & 100 & 100 & 100\\
Cotton & 7.7 & 5.2 & 100 & 100 & 100\\
Tomato & 2.5 & 9.9 & 100 & 100 & 80\\
Potato & 1.7 & 12 & 100 & 93 & 66\\
Soybean & 5.0 & 20 & 100 & 100 & 100\\
Sugar beet & 7.0 & 5.9 & 100 & 100 & 100\\
Sunflower & 4.8 & 5.0 & 100 & 100 & 100\\
Onion & 1.2 & 16 & 100 & 83 & 47\\
Alfalfa & 2.0 & 7.3 & 100 & 98 & 82\\\bottomrule\end{tabular}\end{table}

\exbox{Salinity under drip}{Onion ($a=1.2$, $b=16$) and cotton ($a=7.7$, $b=5.2$) are drip-irrigated with water of 2.0 dS\,m$^{-1}$. Compare the relative yields and say what the leaching requirement of Chapter~14 would be.}\sol \givenline{$a=1.2$ ; $b=16$ ; $a=7.7$ ; $b=5.2$}\ans $EC_e=1.5\times2.0=3.0$ dS\,m$^{-1}$ \ans Onion: $1-0.16(3.0-1.2)=0.712$, i.e.\ 71.2 per cent \ans Cotton: $3.0<7.7$, so $Y/Y_m=1.000$ \ans Leaching requirement $LR=EC_w/(5EC_e-EC_w)$ with the crop's own threshold: onion $2.0/(6.0-2.0)=0.500$; cotton $2.0/(38.5-2.0)=0.055$ \ans The sensitive crop must be leached nine times harder: for it, drip irrigation of saline water needs frequent irrigation to keep the salt at the bulb edge away from the roots, while cotton needs no special care \finans

}{}
\banner{Drip irrigation design}\label{ch:drip}
\section{Depth per irrigation}
By default the depth is the crop need over the interval capped by what the root zone holds at the trigger: $Z_n=\min\{Z_n^{\rm DSS},\max(4,ET_{\rm pk}\Delta t)\}$ mm, with $ET_{\rm pk}=\mathrm{clip}\{\max_d\frac15\sum_{j=0}^{4}(E_{d+j}+T_{d+j});2.5,9\}$ mm d$^{-1}$.

\exbox{Depth per irrigation}{The five-day mean of $E+T$ peaks at $6.4$ mm d$^{-1}$ in the season of the loam. Find $ET_{\rm pk}$, and $Z_n$ for irrigation every 1, 2, 3 and 5 days with $Z_n^{\rm DSS}=73.3$ mm.}\sol  $ET_{\rm pk}=\mathrm{clip}(6.4;2.5,9)=6.4$. $Z_n=\min\{73.3,\max(4,6.4\Delta t)\}$: $\Delta t=1$: $6.4$ mm; $2$: $12.8$; $3$: $19.2$; $5$: $32.0$ mm. All are below the root-zone limit 73.3 mm, so the crop need decides. At the default $\Delta t=2$ d the depth is 12.8 mm, which is what a drip line must deliver at each irrigation.


\section{Lateral hydraulics}
A lateral of length $L_L$ has $N=\mathrm{round}(L_L/S_e)$ intervals and $N+1$ emitters. The friction loss in the segment feeding emitter $j$ is the Hazen--Williams loss $h_{f,j}=\dfrac{10.67\,S_eQ_j^{1.852}}{C^{1.852}D^{4.871}}$ with $C=140$ and $Q_j=\sum_{i\ge j+1}q_i$ (m$^3$ s$^{-1}$); the head is $h_j=h_{j-1}-h_{f,j-1}$ with $h_0=H_0$; the emitter law is $q_j=q(h_j/H_0)^x$ (non-compensating), or for the pressure-compensating emitter $q_j=q[1+0.04(h_j-H_0)/H_0]$ above $H_{\min}=8$ m and $q(\max(h_j,0.1)/H_{\min})^{0.5}$ below it. Flows and heads depend on each other: the pair is iterated from the uniform start (three plain substitution passes, then the flows averaged with the previous pass) until the largest relative change is below $10^{-7}$.

\exbox{Hand pass of a short lateral}{Hand-solve a lateral with three emitters ($N=2$): online emitter $q=4$ L h$^{-1}$ at $H_0=10$ m, $x=0.5$, $S_e=0.75$ m, $D=16$ mm, $C=140$. Do the first pass of the iteration and the check of the loss.}\sol  Start with $q_j=4$ L h$^{-1}$ $=1.111\times10^{-6}$ m$^3$ s$^{-1}$ at every emitter. The segment before emitter 1 carries $Q=q_1+q_2=2.222\times10^{-6}$ m$^3$ s$^{-1}$: $h_f=10.67(0.75)(2.222\times10^{-6})^{1.852}/(140^{1.852}\times0.016^{4.871})$. Compute: $(2.222\times10^{-6})^{1.852}=e^{1.852\ln(2.222\times10^{-6})}=e^{1.852\times(-13.017)}=e^{-24.11}=3.4\times10^{-11}$; $140^{1.852}=e^{1.852\times4.9416}=e^{9.152}=9{,}420$; $0.016^{4.871}=e^{4.871\times(-4.1352)}=e^{-20.14}=1.8\times10^{-9}$. So $h_f=8.0\times3.4\times10^{-11}/(9420\times1.8\times10^{-9})=2.7\times10^{-10}/1.7\times10^{-5}=1.6\times10^{-5}$ m: a negligible loss, so $q_1\approx q_0$ and the three emitters are the same at this length. Short laterals are uniform; the loss grows with the square of the length (flow) times length (pipe) so that it is of order $L^{2.852}$.


\exbox{Flow and head along a 100 m lateral}{Solve a lateral of $L_L=100$ m with the online emitter ($q=4$ L h$^{-1}$, $S_e=0.75$ m, $x=0.5$, $H_0=10$ m, $D=16$ mm): the head at the far end, $\bar q$, $q_{\min}$ and the number of passes.}\sol  $N=\mathrm{round}(100/0.75)=133$ intervals and $134$ emitters. At the start the segment at the inlet carries all 134 emitters: $Q=134\times1.111\times10^{-6}=1.489\times10^{-4}$ m$^3$ s$^{-1}$, and the loss of its first 0.75 m is $h_f=0.03877$ m. Summing the losses along the lateral lowers the head, which lowers the flows, which lowers the losses: Table~\ref{tab:latiter} lists the first ten passes (the first three are plain substitution). The iteration converges (to $10^{-7}$) after 13 passes to $\bar q=3.7590$ L h$^{-1}$, $q_{\min}=3.6722$ L h$^{-1}$ at the last emitter and an end head $H_{\rm end}=8.428$ m. The loss along the lateral is $h_f=H_0-H_{\rm end}=1.57$ m, 15.7 per cent of the inlet head (the warning level is 20 per cent).\begin{table}[!htbp]\centering\small\caption{Convergence of the lateral solution.}\label{tab:latiter}
\begin{tabular}{lrrrr}\toprule
\hd{Pass} & \hd{$\bar q$ (L/h)} & \hd{$q_{\min}$ (L/h)} & \hd{$H_{\rm end}$ (m)} & \hd{max change}\\\midrule
1 & 3.7230 & 3.6217 & 8.198 & 9.46e-02\\
2 & 3.7642 & 3.6795 & 8.462 & 1.60e-02\\
3 & 3.7582 & 3.6711 & 8.423 & 2.27e-03\\
4 & 3.7587 & 3.6717 & 8.429 & 1.65e-04\\
5 & 3.7588 & 3.6720 & 8.428 & 7.06e-05\\
6 & 3.7589 & 3.6721 & 8.428 & 3.02e-05\\
7 & 3.7590 & 3.6722 & 8.428 & 1.29e-05\\
8 & 3.7590 & 3.6722 & 8.428 & 5.51e-06\\
9 & 3.7590 & 3.6722 & 8.428 & 2.36e-06\\
10 & 3.7590 & 3.6722 & 8.428 & 1.01e-06\\\bottomrule\end{tabular}\end{table}


\section{Emission uniformity}
\begin{equation}E_u=\mathrm{clip}\Bigl\{\Bigl(1-\frac{1.27c_v}{\sqrt{n_e}}\Bigr)\frac{q_{\min}}{\bar q};0.3,1\Bigr\},\qquad\sigma_q=\sqrt{c_v^2+\sigma_P^2},\ \sigma_P^2=\frac1{N+1}\sum_j\Bigl(\frac{q_j}{\bar q}-1\Bigr)^2.\end{equation}
\par\noindent where $E_u$ is the emission uniformity (0.3 to 1), $c_v$ the manufacturing coefficient of variation, $n_e$ the number of emitters per plant, $q_{\min}$ and $\bar q$ the minimum and mean discharge (L\,h$^{-1}$), $\sigma_q$ the total variation of discharge, $\sigma_P$ that caused by pressure, $q_j$ the discharge of emitter $j$ and $N+1$ the number of emitters.


\exbox{Emission uniformity of four emitters}{Compute the emission uniformity of the four emitter types of the application (Table 12.1 of the application: online, inline, drip tape and pressure-compensating) on a lateral of 100 m, with $D=16$ mm.}\sol  Solve each lateral as above, then apply the clip formula. For the online emitter: $(1-1.27(0.07)/\sqrt1)=0.911$, $q_{\min}/\bar q=3.672/3.759=0.977$, so $E_u=0.911\times0.977=0.890$, $89$ per cent. Table~\ref{tab:eu} gives the four.\begin{table}[!htbp]\centering\small\caption{Uniformity of four emitters on a 100 m lateral.}\label{tab:eu}
\begin{tabular}{lrrrrrr}\toprule
\hd{Emitter} & \hd{$q$ (L/h)} & \hd{$S_e$ (m)} & \hd{Head loss (\%)} & \hd{$q_{\min}/\bar q$} & \hd{$E_u$ (\%)} & \hd{$\sigma_q$}\\\midrule
Online & 4 & 0.75 & 15.7 & 0.977 & 89.0 & 0.074\\
Inline & 2 & 0.5 & 9.8 & 0.986 & 93.6 & 0.043\\
Drip tape & 1 & 0.2 & 13.8 & 0.976 & 90.6 & 0.084\\
PC & 2 & 0.6 & 5.1 & 0.999 & 96.1 & 0.030\\\bottomrule\end{tabular}\end{table}
\textbf{Reading.} The drip tape (many small emitters: $2\times$ per plant and $S_e=0.2$ m) flows 1 L h$^{-1}$ each, but the 500 emitters on 100 m make the lateral carry 0.14 L s$^{-1}$; its head loss ($13.8$ per cent) is comparable to that of the online emitter. The manufacturing term is the larger part of the non-uniformity for the compensating emitter, whose hydraulic variation is tiny.


\exbox{How long can the lateral be}{Find how the head loss and the flow ratio $q_{\min}/\bar q$ of the online emitter (D 16 mm) and the pressure-compensating emitter change with the lateral length 50 to 300 m, and where the 20 per cent loss warning appears.}\sol  Table~\ref{tab:long} lists the solutions. For the online emitter the loss rises from a few per cent at 50 m to tens of per cent at 300 m, since the loss goes as $L^{2.852}$ in a uniform pipe; the flow ratio falls by roughly half that percentage, because $q\propto h^{0.5}$ (a 10 per cent head loss costs about 5 per cent of flow). The compensating emitter keeps $q_{\min}/\bar q$ near 1 until the head falls below $H_{\min}=8$ m. (Beyond about 250 m the head reaches the application's floor of 0.5 m, which is why the loss column stops growing; such a lateral is outside the design rule.)\begin{table}[!htbp]\centering\small\caption{Effect of lateral length.}\label{tab:long}
\begin{tabular}{lrrrr}\toprule
\hd{$L_L$ (m)} & \hd{online loss (\%)} & \hd{online $q_{\min}/\bar q$} & \hd{PC loss (\%)} & \hd{PC $q_{\min}/\bar q$}\\\midrule
50 & 2.5 & 0.997 & 0.7 & 1.000\\
100 & 15.7 & 0.977 & 5.1 & 0.999\\
150 & 39.0 & 0.930 & 15.9 & 0.998\\
200 & 62.1 & 0.855 & 35.4 & 0.996\\
250 & 78.5 & 0.758 & 59.1 & 0.930\\
300 & 72.9 & 0.751 & 73.8 & 0.837\\\bottomrule\end{tabular}\end{table}


\section{The wetted bulb}
The soil stores $\Delta\theta=\max\{0.03,f_T(\theta_{fc}-\theta_{wp})\}$. The bulb below the emitter is a half-ellipsoid of horizontal radius $r$ and lower semi-axis $ar$, with a short upper half of semi-axis $ur$, $u=0.6$: $\rho(z)=r\sqrt{1-\zeta^2}$ with $\zeta=(z-z_d)/(ar)$ below the emitter and $(z_d-z)/(ur)$ above it. The aspect ratio follows the soil,
\begin{equation}\boxed{a=\mathrm{clip}\Bigl\{\Bigl[0.7+1.6\Bigl(\frac{K_s}{K_s+15}\Bigr)^{1.1}\Bigr]\sqrt{\mathrm{clip}(2/q;0.7,1.4)};0.75,2.4\Bigr\},}\end{equation}
\par\noindent where $a$ is the radius of the wetted bulb (m), $K_s$ the saturated conductivity (cm\,d$^{-1}$) and $q$ the emitter discharge (L\,h$^{-1}$); the constants set the soil and discharge effects, and the value is kept between 0.75 and 2.4 m.

and the radius is found by bisection so that $\int\pi\rho(z)^2dz=V\cdot10^6/\Delta\theta$ mm$^3$ ($V$ in litres).

\exbox{The shape and size of the bulb}{(a) Compute the shape factor $a$ for the ten soils at $q=4$ L h$^{-1}$ (the application quotes about 2.1 for sand, 1.9 sandy loam, 1.3 to 1.4 for loam and silt loam, 0.9 to 1.0 for clay loam and silty clay loam). (b) For the loam with a buried emitter at $z_d=250$ mm delivering 1 litre, find the wetted radius.}
\textbf{Solution (a).} For the loam: $K_s/(K_s+15)=12.07/27.07=0.446$; $0.446^{1.1}=0.4115$; $0.7+1.6\times0.4115=1.358$; $\sqrt{\mathrm{clip}(2/4;0.7,1.4)}=\sqrt{0.7}=0.8367$; $a=1.358\times0.8367=1.136$. For all soils see Table~\ref{tab:aspect}; with $q=4$ L h$^{-1}$ the factor $\sqrt{0.7}$ lowers $a$ by 16 per cent relative to $q=2$, so the values here are lower than the typical figures of the application, which correspond to 2 L h$^{-1}$ ($\sqrt{1}=1$). A larger emitter flow spreads the bulb sideways.

\subhd{Solution (b).} $\Delta\theta=\max\{0.03,0.468(0.2687-0.1309)\}=0.0645$. The wetted volume is $V_w=1\times10^6/0.0645=1.551e+07$ mm$^3$ (about 15.5 litres of soil). Bisection on $r$ in $[5,3500]$ mm over 38 steps until the volume of revolution equals $V_w$ gives $r=162$ mm (lower semi-axis $ar=184$ mm, upper $0.6r=97$ mm). The bulb is buried at 250 mm, so the upper part reaches 153 mm (below the surface).\begin{table}[!htbp]\centering\small\caption{Wetting-bulb aspect ratio from the soil.}\label{tab:aspect}
\begin{tabular}{lrrrr}\toprule
\hd{ID} & \hd{Soil} & \hd{$K_s$ (mm/h)} & \hd{$\Delta\theta$} & \hd{$a$ at $q=4$}\\\midrule
S01 & Coarse loam & 13.1 & 0.057 & 1.17\\
S02 & Coarse sand & 114.7 & 0.030 & 1.76\\
S03 & Sand & 134.7 & 0.030 & 1.78\\
S04 & Loamy sand & 86.4 & 0.030 & 1.71\\
S05 & Sandy loam & 44.7 & 0.042 & 1.56\\
S06 & Loam & 12.1 & 0.062 & 1.14\\
S07 & Silt loam & 10.9 & 0.085 & 1.10\\
S08 & Sandy clay loam & 5.3 & 0.047 & 0.89\\
S09 & Clay loam & 3.2 & 0.062 & 0.78\\
S10 & Silty clay loam & 3.8 & 0.074 & 0.81\\\bottomrule\end{tabular}\end{table}


\exbox{Efficiencies over the spread of emitters}{Evaluate the efficiencies of a drip system on the loam: online emitter ($q=4$ L h$^{-1}$, $\bar q=0.94q=3.76$), $S_e=0.75$ m, laterals $S_l=1.0$ m (cereal), buried at 250 mm, $Z_n=12.8$ mm, root zone $R_{\max}=920$ mm, $E_0=3$ mm d$^{-1}$, $E_u=0.89$, $\sigma_q=\sqrt{0.07^2+0.03^2}=0.076$ (assumed $\sigma_P=0.03$).}\sol  1. Water required per emitter: $V_{req}=Z_nS_eS_l=12.8\times0.75\times1.0=9.60$ L (1 mm over 1 m$^2$ is 1 L). 2. Automatic run time: $t=V_{req}/(\bar qE_u)=9.60/(3.76\times0.89)=2.87$ h. 3. Mean delivery $\bar V=\bar qt=10.79$ L. 4. Five emitters at the normal quantiles $z=-1.28,-0.52,0,0.52,1.28$ deliver $V_i=\bar V(1+\sigma_qz_i)$ (Table~\ref{tab:quant}); for each the bulb radius, the part below $R_{\max}$, the evaporation $V_{ev}=\min\{0.5V,\pi\rho(0)^2\times10^{-6}\times2E_0\}$ and the root water follow Equations (12.10) to (12.15). 5. Storage $V_{st,i}=\min(V_{root,i},V_{req})$ and the excess $V_{ex,i}$. 6. Averages over the five: $E_a=0.890$, $E_s=1.000$, $DP=0.110$, $EV=0.000$; $E_a+DP+EV=1.000$. Gross depth $Z_{gross}=\bar V/(S_eS_l)=14.4$ mm.\begin{table}[!htbp]\centering\small\caption{Five representative emitters.}\label{tab:quant}
\begin{tabular}{lrrrrrrr}\toprule
\hd{$z$} & \hd{$V$ (L)} & \hd{$r$ (mm)} & \hd{$V_{dp}$} & \hd{$V_{ev}$} & \hd{$V_{root}$} & \hd{$V_{st}$} & \hd{$V_{ex}$}\\\midrule
-1.28 & 9.74 & 346 & 0.000 & 0.000 & 9.74 & 9.60 & 0.14\\
-0.52 & 10.36 & 353 & 0.000 & 0.000 & 10.36 & 9.60 & 0.76\\
0 & 10.79 & 358 & 0.000 & 0.000 & 10.79 & 9.60 & 1.19\\
0.52 & 11.21 & 363 & 0.000 & 0.000 & 11.21 & 9.60 & 1.61\\
1.28 & 11.84 & 370 & 0.000 & 0.000 & 11.84 & 9.60 & 2.24\\\bottomrule\end{tabular}\end{table}
\textbf{Reading.} The drier emitters ($z<0$) deliver less than the need, so $E_s<1$; the wetter ones deliver more than the root zone accepts, and the excess is counted with deep percolation. This is why $E_a<1$ even with no water below the root zone: the spread of emitters alone costs efficiency.


\exbox{Soil and placement}{Compare the same drip system on the loam and on the sand ($f_T=0.842$) and with the emitter at the surface, 25 cm and 55 cm.}\sol  Recompute the bulb for each case (Table~\ref{tab:place}). With a surface line the half-spherical bulb exposes the largest wet patch and evaporation takes about 30 per cent of the water; the buried lines lose none. For the three buried depths the bulb stays inside the 92 cm root zone, so the only loss is the excess delivered by the wetter emitters (DP $=0.11$), and the result is the same in the two soils at the same delivery and run time; the soil changes the bulb shape, not the efficiency, for this small depth.\begin{table}[!htbp]\centering\small\caption{Efficiency by soil and placement.}\label{tab:place}
\begin{tabular}{lrrrrrr}\toprule
\hd{Soil} & \hd{Emitter depth} & \hd{$t$ (h)} & \hd{$E_a$} & \hd{$E_s$} & \hd{DP} & \hd{EV}\\\midrule
Loam & 25 cm & 2.87 & 0.890 & 1.000 & 0.110 & 0.000\\
Loam & surface & 2.87 & 0.703 & 0.789 & 0.000 & 0.297\\
Loam & 55 cm & 2.87 & 0.890 & 1.000 & 0.110 & 0.000\\
Sand & 25 cm & 2.87 & 0.890 & 1.000 & 0.110 & 0.000\\
Sand & surface & 2.87 & 0.675 & 0.759 & 0.000 & 0.325\\\bottomrule\end{tabular}\end{table}


\exbox{Layout, pipes and pump of one hectare}{Lay out 1 ha of the loam with laterals of $L_L=100$ m, $S_l=1.0$ m, farm supply $Q_s=30$ L s$^{-1}$, mainline $L_m=120$ m, lift 12 m, inlet head $H_0=10$ m. Find the number of laterals, zones, the zone flow, the pipe diameter, the head losses, the total dynamic head and the pump power.}\sol  $W_f=10^4A_f/L_L=10^4/100=100$ m; $N_L=\lfloor W_f/S_l\rfloor=100$ laterals. Each carries $134\bar q/3600=0.1399$ L s$^{-1}$ ($\bar q=3.759$ L h$^{-1}$), so all together 13.99 L s$^{-1}$. Zones: $n_{zone}=\lceil N_LQ_{lat}/Q_s\rceil=\lceil13.99/30\rceil=1$; laterals per zone $n_l=\lceil100/1\rceil=100$; zone flow $Q_z=n_lQ_{lat}=13.99$ L s$^{-1}$. Pipe: the smallest diameter of the series with $v=4Q_z/(\pi D^2)\le1.5$ m s$^{-1}$ is $D=110$ mm ($v=1.47$ m s$^{-1}$). Submain length $L_{sm}=n_lS_l=100$ m; the submain loss is 0.4 of the full Hazen--Williams loss, $h_{f,sm}=0.78$ m, and the mainline loss $h_{f,m}=2.33$ m. TDH $=z_{lift}+H_0+h_{f,sm}+h_{f,m}+9=12+10+0.78+2.33+9=34.1$ m. Pump power $P=9.81(Q_z/1000)\,TDH/0.65=7.20$ kW. Pipe per hectare: $(N_LL_L+L_{sm})/A_f=10100$ m.


\exbox{Season and farm of the drip system}{Season and farm for the same system with a net seasonal irrigation $y_I=150$ mm and $Z_n=12.8$ mm.}\sol  $n_{ev}=y_I/Z_n=150/12.8=11.7$ events; gross depth $Z_{gross}=14.4$ mm per event; $Z_{season}=n_{ev}Z_{gross}=169$ mm ($=1685$ m$^3$ ha$^{-1}$). Flow per hectare $q_{ha}=q\times10^4/(3600S_eS_l)=4\times10^4/(3600\times0.75\times1.0)=14.81$ L s$^{-1}$ ha$^{-1}$; sets $n_{set}=\lceil q_{ha}A_f/Q_s\rceil=\lceil14.81/30\rceil=1$; round time $t_{round}=1\times2.87=2.87$ h; pumping hours $H_{pump}=n_{ev}t_{round}=34$ h; energy $E_{kWh}=34\times7.20=242$ kWh.
\IfFileExists{extra/proj12.tex}{\banner{Design project 12: a drip lateral and irrigation set for a vegetable block}\label{sec:proj12}
\textbf{Brief.} A lateral of 100 m carries 200 in-line emitters at 0.5-m spacing, with an inlet head of 10 m and a coefficient $x=0.5$. The designer chooses the diameter from the table; peak crop water use is 6 mm\,d$^{-1}$ with a ground cover factor of 0.8, plants at 0.5 by 1.5 m, with one emitter per plant of 1.6 L\,h$^{-1}$.

\begin{table}[!htbp]\centering\small\caption{Lateral diameter alternatives for $L=100$ m with 200 emitters of 1.6 L/h at 10 m head.}
\begin{tabular}{rrrrr}\toprule
\hd{Lateral $D$ (mm)} & \hd{$h_f$ (m)} & \hd{$h_f/H_0$ (\%)} & \hd{$q_{min}$ (L/h)} & \hd{$EU$ (\%)}\\\midrule
14 & 1.18 & 11.8 & 1.502 & 93.2\\
16 & 0.62 & 6.2 & 1.550 & 94.7\\
18 & 0.35 & 3.5 & 1.572 & 95.3\\
20 & 0.21 & 2.1 & 1.583 & 95.7\\\bottomrule\end{tabular}\end{table}

\textbf{Set time.} The daily volume per plant is $4.8\times0.75=3.60$ L; one emitter at 1.6 L\,h$^{-1}$ needs $T=2.2$ h of operation per day. Splitting the lateral into two sets of 100 emitters halves the flow at the head, which cuts the pipe size or the loss. The choice of the designer is the smallest diameter with $EU\ge90$ per cent.
\fig{g12_lat}{0.95}{A 100-m drip lateral with 200 emitters. (a) Pressure head along the lateral for four diameters. (b) The resulting emitter discharge.}
}{}
\IfFileExists{extra/ten12.tex}{\section{Worked solutions on the ten soils}
The wetted radius, and with it the spacing of emitters, depends on the soil.

\exbox{Ten-soil example 12a: Coarse loam (S01)}{Coarse loam ($K_s=31.2$ cm\,d$^{-1}$): an emitter delivers $q=2.0$ L\,h$^{-1}$. Find the wetted radius from the empirical relation of the application, the emitter spacing that overlaps the bulbs, and the number of emitters per hectare at a lateral spacing of 1.5 m.}\sol $K_s/(K_s+15)=0.675$; $0.7+1.6\times0.675^{1.1}=1.739$ \ans $\sqrt{\mathrm{clip}(2/q)}=1.000$, so $a=1.74$ m (limits 0.75 to 2.4 m) \ans Spacing for overlap $S_e=1.6a=2.78$ m, taken as 1.50 m \ans Emitters per hectare: $10^4/(1.5\times1.50)=4444$ \ans

\exbox{Ten-soil example 12a: Coarse sand (S02)}{Coarse sand ($K_s=288.0$ cm\,d$^{-1}$): an emitter delivers $q=2.0$ L\,h$^{-1}$. Find the wetted radius from the empirical relation of the application, the emitter spacing that overlaps the bulbs, and the number of emitters per hectare at a lateral spacing of 1.5 m.}\sol $K_s/(K_s+15)=0.950$; $0.7+1.6\times0.950^{1.1}=2.213$ \ans $\sqrt{\mathrm{clip}(2/q)}=1.000$, so $a=2.21$ m (limits 0.75 to 2.4 m) \ans Spacing for overlap $S_e=1.6a=3.54$ m, taken as 1.50 m \ans Emitters per hectare: $10^4/(1.5\times1.50)=4444$ \ans

\exbox{Ten-soil example 12a: Sand (S03)}{Sand ($K_s=504.0$ cm\,d$^{-1}$): an emitter delivers $q=2.0$ L\,h$^{-1}$. Find the wetted radius from the empirical relation of the application, the emitter spacing that overlaps the bulbs, and the number of emitters per hectare at a lateral spacing of 1.5 m.}\sol $K_s/(K_s+15)=0.971$; $0.7+1.6\times0.971^{1.1}=2.249$ \ans $\sqrt{\mathrm{clip}(2/q)}=1.000$, so $a=2.25$ m (limits 0.75 to 2.4 m) \ans Spacing for overlap $S_e=1.6a=3.60$ m, taken as 1.50 m \ans Emitters per hectare: $10^4/(1.5\times1.50)=4444$ \ans

\exbox{Ten-soil example 12a: Loamy sand (S04)}{Loamy sand ($K_s=146.4$ cm\,d$^{-1}$): an emitter delivers $q=2.0$ L\,h$^{-1}$. Find the wetted radius from the empirical relation of the application, the emitter spacing that overlaps the bulbs, and the number of emitters per hectare at a lateral spacing of 1.5 m.}\sol $K_s/(K_s+15)=0.907$; $0.7+1.6\times0.907^{1.1}=2.137$ \ans $\sqrt{\mathrm{clip}(2/q)}=1.000$, so $a=2.14$ m (limits 0.75 to 2.4 m) \ans Spacing for overlap $S_e=1.6a=3.42$ m, taken as 1.50 m \ans Emitters per hectare: $10^4/(1.5\times1.50)=4444$ \ans

\exbox{Ten-soil example 12a: Sandy loam (S05)}{Sandy loam ($K_s=62.4$ cm\,d$^{-1}$): an emitter delivers $q=2.0$ L\,h$^{-1}$. Find the wetted radius from the empirical relation of the application, the emitter spacing that overlaps the bulbs, and the number of emitters per hectare at a lateral spacing of 1.5 m.}\sol $K_s/(K_s+15)=0.806$; $0.7+1.6\times0.806^{1.1}=1.962$ \ans $\sqrt{\mathrm{clip}(2/q)}=1.000$, so $a=1.96$ m (limits 0.75 to 2.4 m) \ans Spacing for overlap $S_e=1.6a=3.14$ m, taken as 1.50 m \ans Emitters per hectare: $10^4/(1.5\times1.50)=4444$ \ans

\exbox{Ten-soil example 12a: Loam (S06)}{Loam ($K_s=31.2$ cm\,d$^{-1}$): an emitter delivers $q=2.0$ L\,h$^{-1}$. Find the wetted radius from the empirical relation of the application, the emitter spacing that overlaps the bulbs, and the number of emitters per hectare at a lateral spacing of 1.5 m.}\sol $K_s/(K_s+15)=0.675$; $0.7+1.6\times0.675^{1.1}=1.739$ \ans $\sqrt{\mathrm{clip}(2/q)}=1.000$, so $a=1.74$ m (limits 0.75 to 2.4 m) \ans Spacing for overlap $S_e=1.6a=2.78$ m, taken as 1.50 m \ans Emitters per hectare: $10^4/(1.5\times1.50)=4444$ \ans

\exbox{Ten-soil example 12a: Silt loam (S07)}{Silt loam ($K_s=16.3$ cm\,d$^{-1}$): an emitter delivers $q=2.0$ L\,h$^{-1}$. Find the wetted radius from the empirical relation of the application, the emitter spacing that overlaps the bulbs, and the number of emitters per hectare at a lateral spacing of 1.5 m.}\sol $K_s/(K_s+15)=0.521$; $0.7+1.6\times0.521^{1.1}=1.481$ \ans $\sqrt{\mathrm{clip}(2/q)}=1.000$, so $a=1.48$ m (limits 0.75 to 2.4 m) \ans Spacing for overlap $S_e=1.6a=2.37$ m, taken as 1.50 m \ans Emitters per hectare: $10^4/(1.5\times1.50)=4444$ \ans

\exbox{Ten-soil example 12a: Sandy clay loam (S08)}{Sandy clay loam ($K_s=10.3$ cm\,d$^{-1}$): an emitter delivers $q=2.0$ L\,h$^{-1}$. Find the wetted radius from the empirical relation of the application, the emitter spacing that overlaps the bulbs, and the number of emitters per hectare at a lateral spacing of 1.5 m.}\sol $K_s/(K_s+15)=0.408$; $0.7+1.6\times0.408^{1.1}=1.296$ \ans $\sqrt{\mathrm{clip}(2/q)}=1.000$, so $a=1.30$ m (limits 0.75 to 2.4 m) \ans Spacing for overlap $S_e=1.6a=2.07$ m, taken as 1.50 m \ans Emitters per hectare: $10^4/(1.5\times1.50)=4444$ \ans

\exbox{Ten-soil example 12a: Clay loam (S09)}{Clay loam ($K_s=5.5$ cm\,d$^{-1}$): an emitter delivers $q=2.0$ L\,h$^{-1}$. Find the wetted radius from the empirical relation of the application, the emitter spacing that overlaps the bulbs, and the number of emitters per hectare at a lateral spacing of 1.5 m.}\sol $K_s/(K_s+15)=0.269$; $0.7+1.6\times0.269^{1.1}=1.077$ \ans $\sqrt{\mathrm{clip}(2/q)}=1.000$, so $a=1.08$ m (limits 0.75 to 2.4 m) \ans Spacing for overlap $S_e=1.6a=1.72$ m, taken as 1.50 m \ans Emitters per hectare: $10^4/(1.5\times1.50)=4444$ \ans

\exbox{Ten-soil example 12a: Silty clay loam (S10)}{Silty clay loam ($K_s=3.6$ cm\,d$^{-1}$): an emitter delivers $q=2.0$ L\,h$^{-1}$. Find the wetted radius from the empirical relation of the application, the emitter spacing that overlaps the bulbs, and the number of emitters per hectare at a lateral spacing of 1.5 m.}\sol $K_s/(K_s+15)=0.194$; $0.7+1.6\times0.194^{1.1}=0.963$ \ans $\sqrt{\mathrm{clip}(2/q)}=1.000$, so $a=0.96$ m (limits 0.75 to 2.4 m) \ans Spacing for overlap $S_e=1.6a=1.54$ m, taken as 1.50 m \ans Emitters per hectare: $10^4/(1.5\times1.50)=4444$ \ans

}{}
\IfFileExists{extra/tenb12.tex}{\section{Worked solutions on ten crops}
The ten crops of the library differ in their demand, their roots and their sensitivity. The same type of question is answered for each of them at $ET_0=5.5$ mm\,d$^{-1}$; the data are typical values of the crops (Doorenbos--Kassam and FAO 56).

\exbox{Crop example 12a: Maize}{Maize: drip irrigation with plants at 0.25 m along rows 0.75 m apart, one emitter per plant of 2 L\,h$^{-1}$, $K_r=0.8$ ground cover. Find the plants per hectare, the daily water per plant and the run time.}\sol Plants per hectare $10^4/(0.25\times0.75)=53333$ \ans $V=K_cET_0\,K_r\,S_pS_r=1.20\times5.5\times0.8\times0.25\times0.75=0.99$ L per plant per day \ans Run time $V/q=0.99/2=0.49$ h \ans

\exbox{Crop example 12a: Wheat}{Wheat: drip irrigation with plants at 0.15 m along rows 0.75 m apart, one emitter per plant of 2 L\,h$^{-1}$, $K_r=0.8$ ground cover. Find the plants per hectare, the daily water per plant and the run time.}\sol Plants per hectare $10^4/(0.15\times0.75)=88889$ \ans $V=K_cET_0\,K_r\,S_pS_r=1.15\times5.5\times0.8\times0.15\times0.75=0.57$ L per plant per day \ans Run time $V/q=0.57/2=0.28$ h \ans

\exbox{Crop example 12a: Cotton}{Cotton: drip irrigation with plants at 0.15 m along rows 1.00 m apart, one emitter per plant of 2 L\,h$^{-1}$, $K_r=0.8$ ground cover. Find the plants per hectare, the daily water per plant and the run time.}\sol Plants per hectare $10^4/(0.15\times1.00)=66667$ \ans $V=K_cET_0\,K_r\,S_pS_r=1.15\times5.5\times0.8\times0.15\times1.00=0.76$ L per plant per day \ans Run time $V/q=0.76/2=0.38$ h \ans

\exbox{Crop example 12a: Tomato}{Tomato: drip irrigation with plants at 0.50 m along rows 1.50 m apart, one emitter per plant of 2 L\,h$^{-1}$, $K_r=0.8$ ground cover. Find the plants per hectare, the daily water per plant and the run time.}\sol Plants per hectare $10^4/(0.50\times1.50)=13333$ \ans $V=K_cET_0\,K_r\,S_pS_r=1.15\times5.5\times0.8\times0.50\times1.50=3.79$ L per plant per day \ans Run time $V/q=3.79/2=1.90$ h \ans

\exbox{Crop example 12a: Potato}{Potato: drip irrigation with plants at 0.30 m along rows 0.75 m apart, one emitter per plant of 2 L\,h$^{-1}$, $K_r=0.8$ ground cover. Find the plants per hectare, the daily water per plant and the run time.}\sol Plants per hectare $10^4/(0.30\times0.75)=44444$ \ans $V=K_cET_0\,K_r\,S_pS_r=1.15\times5.5\times0.8\times0.30\times0.75=1.14$ L per plant per day \ans Run time $V/q=1.14/2=0.57$ h \ans

\exbox{Crop example 12a: Soybean}{Soybean: drip irrigation with plants at 0.15 m along rows 0.50 m apart, one emitter per plant of 2 L\,h$^{-1}$, $K_r=0.8$ ground cover. Find the plants per hectare, the daily water per plant and the run time.}\sol Plants per hectare $10^4/(0.15\times0.50)=133333$ \ans $V=K_cET_0\,K_r\,S_pS_r=1.15\times5.5\times0.8\times0.15\times0.50=0.38$ L per plant per day \ans Run time $V/q=0.38/2=0.19$ h \ans

\exbox{Crop example 12a: Sugar beet}{Sugar beet: drip irrigation with plants at 0.20 m along rows 0.50 m apart, one emitter per plant of 2 L\,h$^{-1}$, $K_r=0.8$ ground cover. Find the plants per hectare, the daily water per plant and the run time.}\sol Plants per hectare $10^4/(0.20\times0.50)=100000$ \ans $V=K_cET_0\,K_r\,S_pS_r=1.20\times5.5\times0.8\times0.20\times0.50=0.53$ L per plant per day \ans Run time $V/q=0.53/2=0.26$ h \ans

\exbox{Crop example 12a: Sunflower}{Sunflower: drip irrigation with plants at 0.30 m along rows 0.75 m apart, one emitter per plant of 2 L\,h$^{-1}$, $K_r=0.8$ ground cover. Find the plants per hectare, the daily water per plant and the run time.}\sol Plants per hectare $10^4/(0.30\times0.75)=44444$ \ans $V=K_cET_0\,K_r\,S_pS_r=1.10\times5.5\times0.8\times0.30\times0.75=1.09$ L per plant per day \ans Run time $V/q=1.09/2=0.54$ h \ans

\exbox{Crop example 12a: Onion}{Onion: drip irrigation with plants at 0.15 m along rows 0.40 m apart, one emitter per plant of 2 L\,h$^{-1}$, $K_r=0.8$ ground cover. Find the plants per hectare, the daily water per plant and the run time.}\sol Plants per hectare $10^4/(0.15\times0.40)=166667$ \ans $V=K_cET_0\,K_r\,S_pS_r=1.05\times5.5\times0.8\times0.15\times0.40=0.28$ L per plant per day \ans Run time $V/q=0.28/2=0.14$ h \ans

\exbox{Crop example 12a: Alfalfa}{Alfalfa: drip irrigation with plants at 0.15 m along rows 0.75 m apart, one emitter per plant of 2 L\,h$^{-1}$, $K_r=0.8$ ground cover. Find the plants per hectare, the daily water per plant and the run time.}\sol Plants per hectare $10^4/(0.15\times0.75)=88889$ \ans $V=K_cET_0\,K_r\,S_pS_r=1.15\times5.5\times0.8\times0.15\times0.75=0.57$ L per plant per day \ans Run time $V/q=0.57/2=0.28$ h \ans

}{}
\IfFileExists{extra/teny12.tex}{\section{Worked solutions on the ten soils (third scenario)}
The same questions are repeated for different working conditions, to show how the answer moves with them.

\exbox{Ten-soil example 12c: Coarse loam (S01)}{Coarse loam ($K_s=31.2$ cm\,d$^{-1}$): an emitter delivers $q=4.0$ L\,h$^{-1}$. Find the wetted radius from the empirical relation of the application, the emitter spacing that overlaps the bulbs, and the number of emitters per hectare at a lateral spacing of 1.5 m.}\sol $K_s/(K_s+15)=0.675$; $0.7+1.6\times0.675^{1.1}=1.739$ \ans $\sqrt{\mathrm{clip}(2/q)}=0.837$, so $a=1.45$ m (limits 0.75 to 2.4 m) \ans Spacing for overlap $S_e=1.6a=2.33$ m, taken as 1.50 m \ans Emitters per hectare: $10^4/(1.5\times1.50)=4444$ \ans

\exbox{Ten-soil example 12c: Coarse sand (S02)}{Coarse sand ($K_s=288.0$ cm\,d$^{-1}$): an emitter delivers $q=4.0$ L\,h$^{-1}$. Find the wetted radius from the empirical relation of the application, the emitter spacing that overlaps the bulbs, and the number of emitters per hectare at a lateral spacing of 1.5 m.}\sol $K_s/(K_s+15)=0.950$; $0.7+1.6\times0.950^{1.1}=2.213$ \ans $\sqrt{\mathrm{clip}(2/q)}=0.837$, so $a=1.85$ m (limits 0.75 to 2.4 m) \ans Spacing for overlap $S_e=1.6a=2.96$ m, taken as 1.50 m \ans Emitters per hectare: $10^4/(1.5\times1.50)=4444$ \ans

\exbox{Ten-soil example 12c: Sand (S03)}{Sand ($K_s=504.0$ cm\,d$^{-1}$): an emitter delivers $q=4.0$ L\,h$^{-1}$. Find the wetted radius from the empirical relation of the application, the emitter spacing that overlaps the bulbs, and the number of emitters per hectare at a lateral spacing of 1.5 m.}\sol $K_s/(K_s+15)=0.971$; $0.7+1.6\times0.971^{1.1}=2.249$ \ans $\sqrt{\mathrm{clip}(2/q)}=0.837$, so $a=1.88$ m (limits 0.75 to 2.4 m) \ans Spacing for overlap $S_e=1.6a=3.01$ m, taken as 1.50 m \ans Emitters per hectare: $10^4/(1.5\times1.50)=4444$ \ans

\exbox{Ten-soil example 12c: Loamy sand (S04)}{Loamy sand ($K_s=146.4$ cm\,d$^{-1}$): an emitter delivers $q=4.0$ L\,h$^{-1}$. Find the wetted radius from the empirical relation of the application, the emitter spacing that overlaps the bulbs, and the number of emitters per hectare at a lateral spacing of 1.5 m.}\sol $K_s/(K_s+15)=0.907$; $0.7+1.6\times0.907^{1.1}=2.137$ \ans $\sqrt{\mathrm{clip}(2/q)}=0.837$, so $a=1.79$ m (limits 0.75 to 2.4 m) \ans Spacing for overlap $S_e=1.6a=2.86$ m, taken as 1.50 m \ans Emitters per hectare: $10^4/(1.5\times1.50)=4444$ \ans

\exbox{Ten-soil example 12c: Sandy loam (S05)}{Sandy loam ($K_s=62.4$ cm\,d$^{-1}$): an emitter delivers $q=4.0$ L\,h$^{-1}$. Find the wetted radius from the empirical relation of the application, the emitter spacing that overlaps the bulbs, and the number of emitters per hectare at a lateral spacing of 1.5 m.}\sol $K_s/(K_s+15)=0.806$; $0.7+1.6\times0.806^{1.1}=1.962$ \ans $\sqrt{\mathrm{clip}(2/q)}=0.837$, so $a=1.64$ m (limits 0.75 to 2.4 m) \ans Spacing for overlap $S_e=1.6a=2.63$ m, taken as 1.50 m \ans Emitters per hectare: $10^4/(1.5\times1.50)=4444$ \ans

\exbox{Ten-soil example 12c: Loam (S06)}{Loam ($K_s=31.2$ cm\,d$^{-1}$): an emitter delivers $q=4.0$ L\,h$^{-1}$. Find the wetted radius from the empirical relation of the application, the emitter spacing that overlaps the bulbs, and the number of emitters per hectare at a lateral spacing of 1.5 m.}\sol $K_s/(K_s+15)=0.675$; $0.7+1.6\times0.675^{1.1}=1.739$ \ans $\sqrt{\mathrm{clip}(2/q)}=0.837$, so $a=1.45$ m (limits 0.75 to 2.4 m) \ans Spacing for overlap $S_e=1.6a=2.33$ m, taken as 1.50 m \ans Emitters per hectare: $10^4/(1.5\times1.50)=4444$ \ans

\exbox{Ten-soil example 12c: Silt loam (S07)}{Silt loam ($K_s=16.3$ cm\,d$^{-1}$): an emitter delivers $q=4.0$ L\,h$^{-1}$. Find the wetted radius from the empirical relation of the application, the emitter spacing that overlaps the bulbs, and the number of emitters per hectare at a lateral spacing of 1.5 m.}\sol $K_s/(K_s+15)=0.521$; $0.7+1.6\times0.521^{1.1}=1.481$ \ans $\sqrt{\mathrm{clip}(2/q)}=0.837$, so $a=1.24$ m (limits 0.75 to 2.4 m) \ans Spacing for overlap $S_e=1.6a=1.98$ m, taken as 1.50 m \ans Emitters per hectare: $10^4/(1.5\times1.50)=4444$ \ans

\exbox{Ten-soil example 12c: Sandy clay loam (S08)}{Sandy clay loam ($K_s=10.3$ cm\,d$^{-1}$): an emitter delivers $q=4.0$ L\,h$^{-1}$. Find the wetted radius from the empirical relation of the application, the emitter spacing that overlaps the bulbs, and the number of emitters per hectare at a lateral spacing of 1.5 m.}\sol $K_s/(K_s+15)=0.408$; $0.7+1.6\times0.408^{1.1}=1.296$ \ans $\sqrt{\mathrm{clip}(2/q)}=0.837$, so $a=1.08$ m (limits 0.75 to 2.4 m) \ans Spacing for overlap $S_e=1.6a=1.74$ m, taken as 1.50 m \ans Emitters per hectare: $10^4/(1.5\times1.50)=4444$ \ans

\exbox{Ten-soil example 12c: Clay loam (S09)}{Clay loam ($K_s=5.5$ cm\,d$^{-1}$): an emitter delivers $q=4.0$ L\,h$^{-1}$. Find the wetted radius from the empirical relation of the application, the emitter spacing that overlaps the bulbs, and the number of emitters per hectare at a lateral spacing of 1.5 m.}\sol $K_s/(K_s+15)=0.269$; $0.7+1.6\times0.269^{1.1}=1.077$ \ans $\sqrt{\mathrm{clip}(2/q)}=0.837$, so $a=0.90$ m (limits 0.75 to 2.4 m) \ans Spacing for overlap $S_e=1.6a=1.44$ m, taken as 1.44 m \ans Emitters per hectare: $10^4/(1.5\times1.44)=4622$ \ans

\exbox{Ten-soil example 12c: Silty clay loam (S10)}{Silty clay loam ($K_s=3.6$ cm\,d$^{-1}$): an emitter delivers $q=4.0$ L\,h$^{-1}$. Find the wetted radius from the empirical relation of the application, the emitter spacing that overlaps the bulbs, and the number of emitters per hectare at a lateral spacing of 1.5 m.}\sol $K_s/(K_s+15)=0.194$; $0.7+1.6\times0.194^{1.1}=0.963$ \ans $\sqrt{\mathrm{clip}(2/q)}=0.837$, so $a=0.81$ m (limits 0.75 to 2.4 m) \ans Spacing for overlap $S_e=1.6a=1.29$ m, taken as 1.29 m \ans Emitters per hectare: $10^4/(1.5\times1.29)=5173$ \ans

}{}
\IfFileExists{extra/tenby12.tex}{\subsection{Worked solutions on ten crops (hot climate)}
The ten crops of the library differ in their demand, their roots and their sensitivity. The same type of question is answered for each of them at $ET_0=7.5$ mm\,d$^{-1}$; the data are typical values of the crops (Doorenbos--Kassam and FAO 56).

\exbox{Crop example 12c: Maize}{Maize: drip irrigation with plants at 0.25 m along rows 0.75 m apart, one emitter per plant of 2 L\,h$^{-1}$, $K_r=0.8$ ground cover. Find the plants per hectare, the daily water per plant and the run time.}\sol \givenline{$K_r=0.8$ ground cover}\ans Plants per hectare $10^4/(0.25\times0.75)=53333$ \ans $V=K_cET_0\,K_r\,S_pS_r=1.20\times7.5\times0.8\times0.25\times0.75=1.35$ L per plant per day \ans Run time $V/q=1.35/2=0.68$ h \finans

\exbox{Crop example 12c: Wheat}{Wheat: drip irrigation with plants at 0.15 m along rows 0.75 m apart, one emitter per plant of 2 L\,h$^{-1}$, $K_r=0.8$ ground cover. Find the plants per hectare, the daily water per plant and the run time.}\sol \givenline{$K_r=0.8$ ground cover}\ans Plants per hectare $10^4/(0.15\times0.75)=88889$ \ans $V=K_cET_0\,K_r\,S_pS_r=1.15\times7.5\times0.8\times0.15\times0.75=0.78$ L per plant per day \ans Run time $V/q=0.78/2=0.39$ h \finans

\exbox{Crop example 12c: Cotton}{Cotton: drip irrigation with plants at 0.15 m along rows 1.00 m apart, one emitter per plant of 2 L\,h$^{-1}$, $K_r=0.8$ ground cover. Find the plants per hectare, the daily water per plant and the run time.}\sol \givenline{$K_r=0.8$ ground cover}\ans Plants per hectare $10^4/(0.15\times1.00)=66667$ \ans $V=K_cET_0\,K_r\,S_pS_r=1.15\times7.5\times0.8\times0.15\times1.00=1.03$ L per plant per day \ans Run time $V/q=1.03/2=0.52$ h \finans

\exbox{Crop example 12c: Tomato}{Tomato: drip irrigation with plants at 0.50 m along rows 1.50 m apart, one emitter per plant of 2 L\,h$^{-1}$, $K_r=0.8$ ground cover. Find the plants per hectare, the daily water per plant and the run time.}\sol \givenline{$K_r=0.8$ ground cover}\ans Plants per hectare $10^4/(0.50\times1.50)=13333$ \ans $V=K_cET_0\,K_r\,S_pS_r=1.15\times7.5\times0.8\times0.50\times1.50=5.18$ L per plant per day \ans Run time $V/q=5.18/2=2.59$ h \finans

\exbox{Crop example 12c: Potato}{Potato: drip irrigation with plants at 0.30 m along rows 0.75 m apart, one emitter per plant of 2 L\,h$^{-1}$, $K_r=0.8$ ground cover. Find the plants per hectare, the daily water per plant and the run time.}\sol \givenline{$K_r=0.8$ ground cover}\ans Plants per hectare $10^4/(0.30\times0.75)=44444$ \ans $V=K_cET_0\,K_r\,S_pS_r=1.15\times7.5\times0.8\times0.30\times0.75=1.55$ L per plant per day \ans Run time $V/q=1.55/2=0.78$ h \finans

\exbox{Crop example 12c: Soybean}{Soybean: drip irrigation with plants at 0.15 m along rows 0.50 m apart, one emitter per plant of 2 L\,h$^{-1}$, $K_r=0.8$ ground cover. Find the plants per hectare, the daily water per plant and the run time.}\sol \givenline{$K_r=0.8$ ground cover}\ans Plants per hectare $10^4/(0.15\times0.50)=133333$ \ans $V=K_cET_0\,K_r\,S_pS_r=1.15\times7.5\times0.8\times0.15\times0.50=0.52$ L per plant per day \ans Run time $V/q=0.52/2=0.26$ h \finans

\exbox{Crop example 12c: Sugar beet}{Sugar beet: drip irrigation with plants at 0.20 m along rows 0.50 m apart, one emitter per plant of 2 L\,h$^{-1}$, $K_r=0.8$ ground cover. Find the plants per hectare, the daily water per plant and the run time.}\sol \givenline{$K_r=0.8$ ground cover}\ans Plants per hectare $10^4/(0.20\times0.50)=100000$ \ans $V=K_cET_0\,K_r\,S_pS_r=1.20\times7.5\times0.8\times0.20\times0.50=0.72$ L per plant per day \ans Run time $V/q=0.72/2=0.36$ h \finans

\exbox{Crop example 12c: Sunflower}{Sunflower: drip irrigation with plants at 0.30 m along rows 0.75 m apart, one emitter per plant of 2 L\,h$^{-1}$, $K_r=0.8$ ground cover. Find the plants per hectare, the daily water per plant and the run time.}\sol \givenline{$K_r=0.8$ ground cover}\ans Plants per hectare $10^4/(0.30\times0.75)=44444$ \ans $V=K_cET_0\,K_r\,S_pS_r=1.10\times7.5\times0.8\times0.30\times0.75=1.49$ L per plant per day \ans Run time $V/q=1.49/2=0.74$ h \finans

\exbox{Crop example 12c: Onion}{Onion: drip irrigation with plants at 0.15 m along rows 0.40 m apart, one emitter per plant of 2 L\,h$^{-1}$, $K_r=0.8$ ground cover. Find the plants per hectare, the daily water per plant and the run time.}\sol \givenline{$K_r=0.8$ ground cover}\ans Plants per hectare $10^4/(0.15\times0.40)=166667$ \ans $V=K_cET_0\,K_r\,S_pS_r=1.05\times7.5\times0.8\times0.15\times0.40=0.38$ L per plant per day \ans Run time $V/q=0.38/2=0.19$ h \finans

\exbox{Crop example 12c: Alfalfa}{Alfalfa: drip irrigation with plants at 0.15 m along rows 0.75 m apart, one emitter per plant of 2 L\,h$^{-1}$, $K_r=0.8$ ground cover. Find the plants per hectare, the daily water per plant and the run time.}\sol \givenline{$K_r=0.8$ ground cover}\ans Plants per hectare $10^4/(0.15\times0.75)=88889$ \ans $V=K_cET_0\,K_r\,S_pS_r=1.15\times7.5\times0.8\times0.15\times0.75=0.78$ L per plant per day \ans Run time $V/q=0.78/2=0.39$ h \finans

}{}
\IfFileExists{extra/tenx12.tex}{\subsection{Worked solutions on the ten soils (second scenario)}
The same questions are repeated for different working conditions, to show how the answer moves with them.

\exbox{Ten-soil example 12b: Coarse loam (S01)}{Coarse loam ($K_s=31.2$ cm\,d$^{-1}$): an emitter delivers $q=1.2$ L\,h$^{-1}$. Find the wetted radius from the empirical relation of the application, the emitter spacing that overlaps the bulbs, and the number of emitters per hectare at a lateral spacing of 1.5 m.}\sol \givenline{$K_s=31.2$ cm\,d$^{-1}$ ; $q=1.2$ L\,h$^{-1}$}\ans $K_s/(K_s+15)=0.675$; $0.7+1.6\times0.675^{1.1}=1.739$ \ans $\sqrt{\mathrm{clip}(2/q)}=1.183$, so $a=2.06$ m (limits 0.75 to 2.4 m) \ans Spacing for overlap $S_e=1.6a=3.29$ m, taken as 1.50 m \ans Emitters per hectare: $10^4/(1.5\times1.50)=4444$ \finans

\exbox{Ten-soil example 12b: Coarse sand (S02)}{Coarse sand ($K_s=288.0$ cm\,d$^{-1}$): an emitter delivers $q=1.2$ L\,h$^{-1}$. Find the wetted radius from the empirical relation of the application, the emitter spacing that overlaps the bulbs, and the number of emitters per hectare at a lateral spacing of 1.5 m.}\sol \givenline{$K_s=288.0$ cm\,d$^{-1}$ ; $q=1.2$ L\,h$^{-1}$}\ans $K_s/(K_s+15)=0.950$; $0.7+1.6\times0.950^{1.1}=2.213$ \ans $\sqrt{\mathrm{clip}(2/q)}=1.183$, so $a=2.40$ m (limits 0.75 to 2.4 m) \ans Spacing for overlap $S_e=1.6a=3.84$ m, taken as 1.50 m \ans Emitters per hectare: $10^4/(1.5\times1.50)=4444$ \finans

\exbox{Ten-soil example 12b: Sand (S03)}{Sand ($K_s=504.0$ cm\,d$^{-1}$): an emitter delivers $q=1.2$ L\,h$^{-1}$. Find the wetted radius from the empirical relation of the application, the emitter spacing that overlaps the bulbs, and the number of emitters per hectare at a lateral spacing of 1.5 m.}\sol \givenline{$K_s=504.0$ cm\,d$^{-1}$ ; $q=1.2$ L\,h$^{-1}$}\ans $K_s/(K_s+15)=0.971$; $0.7+1.6\times0.971^{1.1}=2.249$ \ans $\sqrt{\mathrm{clip}(2/q)}=1.183$, so $a=2.40$ m (limits 0.75 to 2.4 m) \ans Spacing for overlap $S_e=1.6a=3.84$ m, taken as 1.50 m \ans Emitters per hectare: $10^4/(1.5\times1.50)=4444$ \finans

\exbox{Ten-soil example 12b: Loamy sand (S04)}{Loamy sand ($K_s=146.4$ cm\,d$^{-1}$): an emitter delivers $q=1.2$ L\,h$^{-1}$. Find the wetted radius from the empirical relation of the application, the emitter spacing that overlaps the bulbs, and the number of emitters per hectare at a lateral spacing of 1.5 m.}\sol \givenline{$K_s=146.4$ cm\,d$^{-1}$ ; $q=1.2$ L\,h$^{-1}$}\ans $K_s/(K_s+15)=0.907$; $0.7+1.6\times0.907^{1.1}=2.137$ \ans $\sqrt{\mathrm{clip}(2/q)}=1.183$, so $a=2.40$ m (limits 0.75 to 2.4 m) \ans Spacing for overlap $S_e=1.6a=3.84$ m, taken as 1.50 m \ans Emitters per hectare: $10^4/(1.5\times1.50)=4444$ \finans

\exbox{Ten-soil example 12b: Sandy loam (S05)}{Sandy loam ($K_s=62.4$ cm\,d$^{-1}$): an emitter delivers $q=1.2$ L\,h$^{-1}$. Find the wetted radius from the empirical relation of the application, the emitter spacing that overlaps the bulbs, and the number of emitters per hectare at a lateral spacing of 1.5 m.}\sol \givenline{$K_s=62.4$ cm\,d$^{-1}$ ; $q=1.2$ L\,h$^{-1}$}\ans $K_s/(K_s+15)=0.806$; $0.7+1.6\times0.806^{1.1}=1.962$ \ans $\sqrt{\mathrm{clip}(2/q)}=1.183$, so $a=2.32$ m (limits 0.75 to 2.4 m) \ans Spacing for overlap $S_e=1.6a=3.72$ m, taken as 1.50 m \ans Emitters per hectare: $10^4/(1.5\times1.50)=4444$ \finans

\exbox{Ten-soil example 12b: Loam (S06)}{Loam ($K_s=31.2$ cm\,d$^{-1}$): an emitter delivers $q=1.2$ L\,h$^{-1}$. Find the wetted radius from the empirical relation of the application, the emitter spacing that overlaps the bulbs, and the number of emitters per hectare at a lateral spacing of 1.5 m.}\sol \givenline{$K_s=31.2$ cm\,d$^{-1}$ ; $q=1.2$ L\,h$^{-1}$}\ans $K_s/(K_s+15)=0.675$; $0.7+1.6\times0.675^{1.1}=1.739$ \ans $\sqrt{\mathrm{clip}(2/q)}=1.183$, so $a=2.06$ m (limits 0.75 to 2.4 m) \ans Spacing for overlap $S_e=1.6a=3.29$ m, taken as 1.50 m \ans Emitters per hectare: $10^4/(1.5\times1.50)=4444$ \finans

\exbox{Ten-soil example 12b: Silt loam (S07)}{Silt loam ($K_s=16.3$ cm\,d$^{-1}$): an emitter delivers $q=1.2$ L\,h$^{-1}$. Find the wetted radius from the empirical relation of the application, the emitter spacing that overlaps the bulbs, and the number of emitters per hectare at a lateral spacing of 1.5 m.}\sol \givenline{$K_s=16.3$ cm\,d$^{-1}$ ; $q=1.2$ L\,h$^{-1}$}\ans $K_s/(K_s+15)=0.521$; $0.7+1.6\times0.521^{1.1}=1.481$ \ans $\sqrt{\mathrm{clip}(2/q)}=1.183$, so $a=1.75$ m (limits 0.75 to 2.4 m) \ans Spacing for overlap $S_e=1.6a=2.80$ m, taken as 1.50 m \ans Emitters per hectare: $10^4/(1.5\times1.50)=4444$ \finans

\exbox{Ten-soil example 12b: Sandy clay loam (S08)}{Sandy clay loam ($K_s=10.3$ cm\,d$^{-1}$): an emitter delivers $q=1.2$ L\,h$^{-1}$. Find the wetted radius from the empirical relation of the application, the emitter spacing that overlaps the bulbs, and the number of emitters per hectare at a lateral spacing of 1.5 m.}\sol \givenline{$K_s=10.3$ cm\,d$^{-1}$ ; $q=1.2$ L\,h$^{-1}$}\ans $K_s/(K_s+15)=0.408$; $0.7+1.6\times0.408^{1.1}=1.296$ \ans $\sqrt{\mathrm{clip}(2/q)}=1.183$, so $a=1.53$ m (limits 0.75 to 2.4 m) \ans Spacing for overlap $S_e=1.6a=2.45$ m, taken as 1.50 m \ans Emitters per hectare: $10^4/(1.5\times1.50)=4444$ \finans

\exbox{Ten-soil example 12b: Clay loam (S09)}{Clay loam ($K_s=5.5$ cm\,d$^{-1}$): an emitter delivers $q=1.2$ L\,h$^{-1}$. Find the wetted radius from the empirical relation of the application, the emitter spacing that overlaps the bulbs, and the number of emitters per hectare at a lateral spacing of 1.5 m.}\sol \givenline{$K_s=5.5$ cm\,d$^{-1}$ ; $q=1.2$ L\,h$^{-1}$}\ans $K_s/(K_s+15)=0.269$; $0.7+1.6\times0.269^{1.1}=1.077$ \ans $\sqrt{\mathrm{clip}(2/q)}=1.183$, so $a=1.27$ m (limits 0.75 to 2.4 m) \ans Spacing for overlap $S_e=1.6a=2.04$ m, taken as 1.50 m \ans Emitters per hectare: $10^4/(1.5\times1.50)=4444$ \finans

\exbox{Ten-soil example 12b: Silty clay loam (S10)}{Silty clay loam ($K_s=3.6$ cm\,d$^{-1}$): an emitter delivers $q=1.2$ L\,h$^{-1}$. Find the wetted radius from the empirical relation of the application, the emitter spacing that overlaps the bulbs, and the number of emitters per hectare at a lateral spacing of 1.5 m.}\sol \givenline{$K_s=3.6$ cm\,d$^{-1}$ ; $q=1.2$ L\,h$^{-1}$}\ans $K_s/(K_s+15)=0.194$; $0.7+1.6\times0.194^{1.1}=0.963$ \ans $\sqrt{\mathrm{clip}(2/q)}=1.183$, so $a=1.14$ m (limits 0.75 to 2.4 m) \ans Spacing for overlap $S_e=1.6a=1.82$ m, taken as 1.50 m \ans Emitters per hectare: $10^4/(1.5\times1.50)=4444$ \finans

}{}
\IfFileExists{extra/tenbx12.tex}{\subsection{Worked solutions on ten crops (cool climate)}
The ten crops of the library differ in their demand, their roots and their sensitivity. The same type of question is answered for each of them at $ET_0=4.0$ mm\,d$^{-1}$; the data are typical values of the crops (Doorenbos--Kassam and FAO 56).

\exbox{Crop example 12b: Maize}{Maize: drip irrigation with plants at 0.25 m along rows 0.75 m apart, one emitter per plant of 2 L\,h$^{-1}$, $K_r=0.8$ ground cover. Find the plants per hectare, the daily water per plant and the run time.}\sol \givenline{$K_r=0.8$ ground cover}\ans Plants per hectare $10^4/(0.25\times0.75)=53333$ \ans $V=K_cET_0\,K_r\,S_pS_r=1.20\times4.0\times0.8\times0.25\times0.75=0.72$ L per plant per day \ans Run time $V/q=0.72/2=0.36$ h \finans

\exbox{Crop example 12b: Wheat}{Wheat: drip irrigation with plants at 0.15 m along rows 0.75 m apart, one emitter per plant of 2 L\,h$^{-1}$, $K_r=0.8$ ground cover. Find the plants per hectare, the daily water per plant and the run time.}\sol \givenline{$K_r=0.8$ ground cover}\ans Plants per hectare $10^4/(0.15\times0.75)=88889$ \ans $V=K_cET_0\,K_r\,S_pS_r=1.15\times4.0\times0.8\times0.15\times0.75=0.41$ L per plant per day \ans Run time $V/q=0.41/2=0.21$ h \finans

\exbox{Crop example 12b: Cotton}{Cotton: drip irrigation with plants at 0.15 m along rows 1.00 m apart, one emitter per plant of 2 L\,h$^{-1}$, $K_r=0.8$ ground cover. Find the plants per hectare, the daily water per plant and the run time.}\sol \givenline{$K_r=0.8$ ground cover}\ans Plants per hectare $10^4/(0.15\times1.00)=66667$ \ans $V=K_cET_0\,K_r\,S_pS_r=1.15\times4.0\times0.8\times0.15\times1.00=0.55$ L per plant per day \ans Run time $V/q=0.55/2=0.28$ h \finans

\exbox{Crop example 12b: Tomato}{Tomato: drip irrigation with plants at 0.50 m along rows 1.50 m apart, one emitter per plant of 2 L\,h$^{-1}$, $K_r=0.8$ ground cover. Find the plants per hectare, the daily water per plant and the run time.}\sol \givenline{$K_r=0.8$ ground cover}\ans Plants per hectare $10^4/(0.50\times1.50)=13333$ \ans $V=K_cET_0\,K_r\,S_pS_r=1.15\times4.0\times0.8\times0.50\times1.50=2.76$ L per plant per day \ans Run time $V/q=2.76/2=1.38$ h \finans

\exbox{Crop example 12b: Potato}{Potato: drip irrigation with plants at 0.30 m along rows 0.75 m apart, one emitter per plant of 2 L\,h$^{-1}$, $K_r=0.8$ ground cover. Find the plants per hectare, the daily water per plant and the run time.}\sol \givenline{$K_r=0.8$ ground cover}\ans Plants per hectare $10^4/(0.30\times0.75)=44444$ \ans $V=K_cET_0\,K_r\,S_pS_r=1.15\times4.0\times0.8\times0.30\times0.75=0.83$ L per plant per day \ans Run time $V/q=0.83/2=0.41$ h \finans

\exbox{Crop example 12b: Soybean}{Soybean: drip irrigation with plants at 0.15 m along rows 0.50 m apart, one emitter per plant of 2 L\,h$^{-1}$, $K_r=0.8$ ground cover. Find the plants per hectare, the daily water per plant and the run time.}\sol \givenline{$K_r=0.8$ ground cover}\ans Plants per hectare $10^4/(0.15\times0.50)=133333$ \ans $V=K_cET_0\,K_r\,S_pS_r=1.15\times4.0\times0.8\times0.15\times0.50=0.28$ L per plant per day \ans Run time $V/q=0.28/2=0.14$ h \finans

\exbox{Crop example 12b: Sugar beet}{Sugar beet: drip irrigation with plants at 0.20 m along rows 0.50 m apart, one emitter per plant of 2 L\,h$^{-1}$, $K_r=0.8$ ground cover. Find the plants per hectare, the daily water per plant and the run time.}\sol \givenline{$K_r=0.8$ ground cover}\ans Plants per hectare $10^4/(0.20\times0.50)=100000$ \ans $V=K_cET_0\,K_r\,S_pS_r=1.20\times4.0\times0.8\times0.20\times0.50=0.38$ L per plant per day \ans Run time $V/q=0.38/2=0.19$ h \finans

\exbox{Crop example 12b: Sunflower}{Sunflower: drip irrigation with plants at 0.30 m along rows 0.75 m apart, one emitter per plant of 2 L\,h$^{-1}$, $K_r=0.8$ ground cover. Find the plants per hectare, the daily water per plant and the run time.}\sol \givenline{$K_r=0.8$ ground cover}\ans Plants per hectare $10^4/(0.30\times0.75)=44444$ \ans $V=K_cET_0\,K_r\,S_pS_r=1.10\times4.0\times0.8\times0.30\times0.75=0.79$ L per plant per day \ans Run time $V/q=0.79/2=0.40$ h \finans

\exbox{Crop example 12b: Onion}{Onion: drip irrigation with plants at 0.15 m along rows 0.40 m apart, one emitter per plant of 2 L\,h$^{-1}$, $K_r=0.8$ ground cover. Find the plants per hectare, the daily water per plant and the run time.}\sol \givenline{$K_r=0.8$ ground cover}\ans Plants per hectare $10^4/(0.15\times0.40)=166667$ \ans $V=K_cET_0\,K_r\,S_pS_r=1.05\times4.0\times0.8\times0.15\times0.40=0.20$ L per plant per day \ans Run time $V/q=0.20/2=0.10$ h \finans

\exbox{Crop example 12b: Alfalfa}{Alfalfa: drip irrigation with plants at 0.15 m along rows 0.75 m apart, one emitter per plant of 2 L\,h$^{-1}$, $K_r=0.8$ ground cover. Find the plants per hectare, the daily water per plant and the run time.}\sol \givenline{$K_r=0.8$ ground cover}\ans Plants per hectare $10^4/(0.15\times0.75)=88889$ \ans $V=K_cET_0\,K_r\,S_pS_r=1.15\times4.0\times0.8\times0.15\times0.75=0.41$ L per plant per day \ans Run time $V/q=0.41/2=0.21$ h \finans

}{}
\IfFileExists{extra/ex12.tex}{\backhead{EXERCISES}{khred}
\begin{exlist}
\item Emitter discharge at 7 m if $q_0=2$ at 10 m, $x=0.5$.\exans{1.67 L/h}
\item Christiansen $F$ for $m=1.852$, $N=100$.\exans{0.356}
\item HW loss for $Q=0.1$ L/s, $C=150$, $D=16$, $L=50$.\exans{1.09 m}
\item EU for $q_{min}=1.5$, $q_{avg}=1.6$, $CV=0.04$, $n_e=1$.\exans{89.0}
\item Hours per day if need is 3.2 L per plant, 2 emitters of 1.6.\exans{1.0 h}
\item Name two kinds of drip emitters.\exans{compensating; non-compensating}
\item Why does a larger $D$ raise EU?\exans{lower head loss}
\item What is the acceptable EU range?\exans{80-90 per cent}
\item What is SDI?\exans{subsurface drip}
\item Daily volume per plant for 5 mm/d, $0.4\times1.2$ m, $K_r=0.7$.\exans{1.68 L}
\end{exlist}
}{}
\IfFileExists{extra/dyk12.tex}{\dykband{DO YOU KNOW ?}
\begin{exlist}
\item List the components of a drip system.
\item What does the exponent $x$ measure?
\item What is the Christiansen factor for?
\item How is EU computed?
\item Why do drip systems need filtration?
\item What is a pressure-compensating emitter?
\item What limits the length of a lateral?
\item How is the set time found?
\end{exlist}
}{}
\IfFileExists{extra/obj12.tex}{\dykband{OBJECTIVE TYPE QUESTIONS}
\begin{exlist}
\item $q=$\opt{$K_dH^x$}{$K/H$}{$H/K$}{$K+H$}
\item $x=0$ means\opt{compensating}{laminar}{leaking}{none}
\item $F$ depends on\opt{$N$ and $m$}{rain}{salt}{none}
\item EU above 90 is\opt{excellent}{poor}{none}{fail}
\item Drip needs\opt{filters}{pumps off}{tanks only}{none}
\item HW uses\opt{$C$}{$n$}{$f$}{$\mu$}
\item SDI is\opt{buried}{sprayed}{flood}{none}
\item Smaller $D$\opt{more loss}{less loss}{none}{rain}
\item Lateral length is limited by\opt{head loss}{rain}{soil}{sun}
\item Set time = volume/\opt{discharge}{rain}{depth}{none}
\end{exlist}
\objkey{1.~(a)\quad 2.~(a)\quad 3.~(a)\quad 4.~(a)\quad 5.~(a)\quad 6.~(a)\quad 7.~(a)\quad 8.~(a)\quad 9.~(a)\quad 10.~(a)}
}{}
